\documentclass[11pt]{amsart}
\usepackage{amsmath,amssymb,amsthm}
\usepackage{graphicx}
\usepackage[hidelinks]{hyperref}

\newtheorem{theorem}{Theorem}

\newcommand{\R}{\mathbb R}
\newcommand{\ellT}{\ell_2}
\newcommand{\BS}{Bhatia--\v{S}emrl}

\title{A Bhatia--\v{S}emrl Operator with Infinite-Dimensional Norming Span}
\author{Agent lane 12}
\date{August 9, 2026}

\begin{document}
\maketitle

\section*{Source and target}

Jayanarayanan--Rajpopat prove that if a nonzero operator $T$ has the
\BS{} property and $\operatorname{span}(M_T)$ is finite dimensional, then
$\|T\|_e<\|T\|$.  In Remark 3.3 of \cite{JR}, they ask whether the
finite-dimensionality hypothesis might follow automatically from the
\BS{} property.

\begin{center}
  \includegraphics[width=0.98\linewidth]{figures/open_question_crop.png}
\end{center}

Here
\[
 M_T=\{x\in S_X:\|Tx\|=\|T\|\},
\]
and $T$ has the \BS{} property when, for every operator $A$ with
$T\perp_B A$, there exists $x\in M_T$ such that $Tx\perp_B Ax$.

\section*{Result}

The answer to the question in Remark 3.3 is negative, even for a scalar-valued
operator on a simple reflexive Banach space.  Take
\[
 X=\R\oplus_\infty\ell_2,
 \qquad f:X\longrightarrow\R,
 \qquad f(a,h)=a.
\]
Then $f$ has the \BS{} property, but
$\operatorname{span}(M_f)=X$ is infinite dimensional.

\section*{Proof intuition}

The unit ball of $\R\oplus_\infty\ell_2$ has two large flat faces on which the
first-coordinate functional attains its norm.  Their linear span contains the
entire $\ell_2$ direction.  At the same time, the dual norm is an
$\ell_1$-sum norm.  Thus orthogonality of $f=(1,0)$ to an arbitrary
$g=(\alpha,k)$ is the elementary inequality $|\alpha|\leq\|k\|_2$.  This
inequality lets us choose a point on the norming face at which $g$ vanishes,
which is exactly the pointwise witness required by the \BS{} property.

\begin{theorem}\label{thm:counterexample}
Let $X=\R\oplus_\infty\ell_2$ and let $f\in X'$ be given by
$f(a,h)=a$.  Then $f$, regarded as an operator from $X$ to $\R$, has the
\BS{} property, while
\[
  \operatorname{span}(M_f)=X.
\]
In particular, the span of $M_f$ is infinite dimensional.
\end{theorem}

\begin{proof}
The norm on $X$ is
\[
  \|(a,h)\|=\max\{|a|,\|h\|_2\},
\]
so $\|f\|=1$ and
\begin{equation}\label{eq:norming-set}
  M_f=\{(\varepsilon,h):\varepsilon\in\{-1,1\},\ \|h\|_2\leq1\}.
\end{equation}
The vector $(1,0)$ lies in $M_f$.  If $\|h\|_2\leq1$, then both $(1,h)$
and $(1,0)$ lie in $M_f$, and
\[
  (0,h)=(1,h)-(1,0).
\]
After scalar multiplication this gives $\{0\}\oplus\ell_2\subseteq
\operatorname{span}(M_f)$, while $(1,0)$ supplies the first coordinate.
Hence $\operatorname{span}(M_f)=X$.

It remains to prove the \BS{} property.  Identify
\[
  X'=\R\oplus_1\ell_2.
\]
Let $g=(\alpha,k)\in X'$ and suppose that $f\perp_B g$.  For every
$\lambda\in\R$,
\begin{equation}\label{eq:dual-norm}
 \|f+\lambda g\|
 =|1+\lambda\alpha|+|\lambda|\,\|k\|_2
 \geq1.
\end{equation}
We claim that
\begin{equation}\label{eq:key-ineq}
  |\alpha|\leq\|k\|_2.
\end{equation}
Indeed, if $\alpha>\|k\|_2$, then for all sufficiently small $t>0$,
\[
 \|f-tg\|=1-t\alpha+t\|k\|_2<1,
\]
contradicting \eqref{eq:dual-norm}.  If
$\alpha<-\|k\|_2$, the same contradiction follows from
$\|f+tg\|<1$ for small $t>0$.  This proves \eqref{eq:key-ineq}.

Suppose first that $k\neq0$, and set
\[
  h=-\frac{\alpha}{\|k\|_2^2}k.
\]
By \eqref{eq:key-ineq}, $\|h\|_2\leq1$.  Thus $x=(1,h)$ belongs to
$M_f$ by \eqref{eq:norming-set}, and
\[
  g(x)=\alpha+\langle k,h\rangle
  =\alpha-\alpha=0.
\]
If $k=0$, then \eqref{eq:key-ineq} gives $\alpha=0$, and we take
$x=(1,0)\in M_f$.  In either case, $f(x)=1$ and $g(x)=0$.  Since a nonzero
real scalar is Birkhoff--James orthogonal precisely to the zero scalar, we
have $f(x)\perp_B g(x)$.  This is exactly the \BS{} property for $f$.
\end{proof}

\section*{Relation to known results}

Kim--Lee proved that every functional on a Banach space has the \BS{} property
if and only if the space is reflexive \cite{KL}.  Since
$\R\oplus_\infty\ell_2$ is reflexive, that theorem independently confirms
the property used above.  The calculation \eqref{eq:dual-norm} gives a direct,
self-contained proof for this example and isolates the flat norming face that
answers Remark 3.3.

The example does \emph{not} answer the source's broader Question 1.9, which
asks whether the \BS{} property always implies $\|T\|_e<\|T\|$.  The range of
$f$ is one dimensional, so $f$ is compact and $\|f\|_e=0<1=\|f\|$.

\section*{Verification and novelty bounds}

The argument uses only the elementary duality
$(\R\oplus_\infty\ell_2)'=\R\oplus_1\ell_2$.  The crucial implication
$f\perp_Bg\Rightarrow|\alpha|\leq\|k\|_2$ was checked in both scalar
directions, and the resulting witness $h$ has norm at most one exactly when
that inequality holds.

The run's registry, solution, attempt, and proof-gap indexes were searched for
arXiv:2512.15208 and the phrases ``Bhatia--\v{S}emrl'',
``$\operatorname{span}(M_T)$'', and ``finite dimensional norm attainment
span''.  No existing packet for this question was found.  Bounded web searches
through August 9, 2026, found the source paper and the Kim--Lee reflexive-space
theorem, but no publication explicitly answering Remark 3.3.  Because the
counterexample is a short consequence of a theorem already cited by the
source, its mathematical validity is high-confidence while its novelty should
be labeled medium-low pending expert review.

\begin{thebibliography}{9}
\bibitem{JR}
C. R. Jayanarayanan and R. R. Rajpopat,
\emph{On Bhatia--\v{S}emrl property, strong subdifferentiability and essential
norm of operators on Banach spaces}, arXiv:2512.15208, 2025.

\bibitem{KL}
S. K. Kim and H. J. Lee,
\emph{The Birkhoff--James orthogonality of operators on infinite dimensional
Banach spaces}, Linear Algebra Appl. 582 (2019), 440--451.
\end{thebibliography}

\end{document}
